\documentclass[11pt,a4paper]{article}

\usepackage[utf8]{inputenc}
\usepackage[T1]{fontenc}
\usepackage[ngerman]{babel}
\usepackage{lmodern}
\usepackage{amsmath,amssymb}
\usepackage{graphicx}
\usepackage{booktabs,tabularx,array}
\usepackage[table]{xcolor}
\usepackage{tikz}
\usetikzlibrary{arrows.meta,positioning,shapes.geometric,calc}
\usepackage[most]{tcolorbox}
\usepackage{enumitem}
\usepackage{microtype}
\usepackage{hyperref}
\usepackage[a4paper,left=18mm,right=18mm,top=15mm,bottom=15mm,headheight=0pt]{geometry}

\definecolor{navy}{HTML}{174A63}
\definecolor{blue}{HTML}{2E7698}
\definecolor{lightblue}{HTML}{EAF3F7}
\definecolor{lightgray}{HTML}{F2F3F4}
\definecolor{midgray}{HTML}{66727A}
\definecolor{green}{HTML}{477A63}
\definecolor{orange}{HTML}{B96A2B}

\hypersetup{colorlinks=true,linkcolor=navy,urlcolor=blue,pdftitle={Sicherungsblatt Aufgabe 3a -- Einen Entscheidungsbaum erstellen}}
\setlength{\parindent}{0pt}
\setlength{\parskip}{4pt}
\setlist[itemize]{leftmargin=5mm,itemsep=1.5pt,topsep=2pt}
\setlist[enumerate]{leftmargin=6mm,itemsep=2pt,topsep=2pt}
\renewcommand{\familydefault}{\sfdefault}

\pagestyle{empty}

\newtcolorbox{definition}[1][]{enhanced,breakable,colback=lightblue,colframe=blue,
  boxrule=0.8pt,arc=1.5mm,left=3mm,right=3mm,top=2mm,bottom=2mm,
  fonttitle=\bfseries,title=Definition,#1}
\newtcolorbox{merke}[1][]{enhanced,breakable,colback=lightgray,colframe=navy,
  boxrule=0.9pt,arc=1.5mm,left=3mm,right=3mm,top=2mm,bottom=2mm,
  fonttitle=\bfseries,title=Merke,#1}
\newtcolorbox{beispiel}[1][]{enhanced,breakable,colback=white,colframe=midgray,
  boxrule=0.7pt,arc=1.5mm,left=3mm,right=3mm,top=2mm,bottom=2mm,
  fonttitle=\bfseries,title=Beispiel,#1}
\newtcolorbox{ueberblick}[1][]{enhanced,breakable,colback=lightblue,colframe=navy,
  boxrule=1.1pt,arc=1.5mm,left=3mm,right=3mm,top=2.5mm,bottom=2.5mm,
  fonttitle=\bfseries,title=Überblick,#1}

\newcommand{\fach}[1]{\textbf{\color{navy}#1}}
\newcommand{\sectionline}[1]{%
  \vspace{1mm}{\Large\bfseries\color{navy}#1}\par\vspace{1mm}%
  {\color{blue}\rule{\linewidth}{0.8pt}}\vspace{1mm}}
\newcommand{\sheettitle}[2]{%
  \begin{tcolorbox}[colback=navy,colframe=navy,arc=2mm,boxrule=0pt,left=5mm,right=5mm,top=2.5mm,bottom=2.5mm]
    {\color{white}\fontsize{20}{24}\selectfont\bfseries #1}\par
    \vspace{0.5mm}{\color{white}\large #2}
  \end{tcolorbox}}
\newcommand{\subline}[1]{\vspace{1mm}{\large\bfseries\color{navy}#1}\par}

\definecolor{fishblue}{HTML}{377ED1}
\definecolor{fishorange}{HTML}{E87931}
\definecolor{fishyellow}{HTML}{F2C94C}
\definecolor{fishred}{HTML}{D84B53}
\definecolor{fishdark}{HTML}{23324A}
\definecolor{peaceful}{HTML}{08705A}
\definecolor{hostile}{HTML}{992D35}

% \fisch{(Position)}{Schuppenfarbe}{Punkte 0/1}{Bauchfarbe}{Flossenfarbe}{ID}{Label f/h}{Fehler 0/1}
\newcommand{\fisch}[8]{%
  \begin{scope}[shift={#1}]
    \filldraw[fill=#5,draw=fishdark,line width=0.5pt,line join=round] (0.46,0) -- (0.74,0.21) -- (0.74,-0.21) -- cycle;
    \filldraw[fill=#5,draw=fishdark,line width=0.5pt,line join=round] (-0.12,0.21) -- (0.08,0.36) -- (0.2,0.19) -- cycle;
    \begin{scope}
      \clip (0,0) ellipse (0.5 and 0.26);
      \fill[#2] (-0.6,-0.3) rectangle (0.6,0.3);
      \fill[#4] (-0.6,-0.3) rectangle (0.6,-0.1);
      \ifnum#3=1
        \fill[fishdark] (-0.12,0.09) circle (0.035) (0.1,0.14) circle (0.035) (0.28,0.05) circle (0.035)
          (0.05,-0.01) circle (0.035) (-0.2,-0.03) circle (0.035);
      \fi
    \end{scope}
    \draw[fishdark,line width=0.7pt] (0,0) ellipse (0.5 and 0.26);
    \filldraw[fill=fishdark,draw=white,line width=0.4pt] (-0.32,0.06) circle (0.055);
    \if#7f
      \node[font=\tiny,text=peaceful,align=center,inner sep=0pt] at (0.12,-0.66) {\textbf{#6}\\friedlich};
    \else
      \node[font=\tiny,text=hostile,align=center,inner sep=0pt] at (0.12,-0.66) {\textbf{#6}\\feindselig};
    \fi
    \ifnum#8=1
      \draw[fishred,line width=0.9pt,dashed,rounded corners=1mm] (-0.6,-0.98) rectangle (0.84,0.5);
      \node[fill=fishred,text=white,font=\tiny\bfseries,rounded corners=0.8mm,inner sep=1pt] at (0.12,0.5) {Fehler};
    \fi
  \end{scope}}

\begin{document}

\sheettitle{Entscheidungsbaum mit Fischdaten erstellen}{Entscheidungsbäume · Sicherungsblatt zu Aufgabe 3a · Informatik 11}

\sectionline{1. Fehlklassifikation und Informationsgewinn}

\begin{definition}[title=Definition: Fehlklassifikation und Informationsgewinn]
Ein Knoten sagt für seine Teilmenge das häufigere Label voraus. Alle Daten mit dem anderen Label sind
\fach{Fehlklassifikationen}. Der \fach{Informationsgewinn} eines Splits ist
\[
  IG=\text{Fehler vor dem Aufteilen}-\text{Gesamtfehler nach dem Aufteilen}.
\]
Gewählt wird das Attribut mit dem größten Informationsgewinn.
\end{definition}

\begin{beispiel}[title=Beispiel: erster Split der 9 Trainingsfische]
\begin{center}
\begin{tikzpicture}[font=\small,
  decision/.style={draw=navy,fill=lightblue,rounded corners=1mm,minimum width=27mm,minimum height=8mm},
  group/.style={rounded corners=2mm,line width=0.9pt}]
  % Alle Trainingsfische ohne Split
  \draw[group,draw=midgray,fill=lightgray] (-7.9,-1.75) rectangle (7.9,1.3);
  \node[anchor=west,font=\small\bfseries,text=navy] at (-7.7,0.98) {Ohne Split: alle 9 Trainingsfische};
  \fisch{(-6.88,0)}{fishorange}{0}{white}{fishyellow}{F1}{h}{0}
  \fisch{(-5.16,0)}{fishorange}{0}{white}{fishred}{F2}{h}{0}
  \fisch{(-3.44,0)}{fishblue}{0}{fishdark}{fishyellow}{F3}{f}{1}
  \fisch{(-1.72,0)}{fishblue}{0}{white}{fishyellow}{F4}{f}{1}
  \fisch{(0,0)}{fishblue}{1}{fishdark}{fishred}{F5}{h}{0}
  \fisch{(1.72,0)}{fishblue}{1}{fishdark}{fishyellow}{F6}{h}{0}
  \fisch{(3.44,0)}{fishorange}{1}{white}{fishyellow}{F7}{h}{0}
  \fisch{(5.16,0)}{fishblue}{1}{white}{fishred}{F8}{f}{1}
  \fisch{(6.88,0)}{fishorange}{1}{fishdark}{fishred}{F9}{f}{1}
  \node[font=\footnotesize] at (0,-1.4) {4 friedlich, 5 feindselig $\;\Rightarrow\;$ Vorhersage \emph{feindselig} $\;\Rightarrow\;$ \textbf{4 Fehler}};

  % Split
  \node[decision] (split) at (0,-2.55) {Schuppenfarbe?};
  \draw (0,-1.75) -- (split.north);
  \draw (split.south west) -- node[above left,font=\footnotesize]{Blau} (-4.0,-3.45);
  \draw (split.south east) -- node[above right,font=\footnotesize]{Orange} (4.0,-3.45);

  % Teilmenge Blau
  \draw[group,draw=fishblue,fill=white] (-7.9,-3.45) rectangle (-0.1,-6.5);
  \node[anchor=west,font=\small\bfseries,text=fishblue] at (-7.7,-3.77) {Teilmenge Blau: 5 Fische};
  \fisch{(-7.05,-4.75)}{fishblue}{0}{fishdark}{fishyellow}{F3}{f}{0}
  \fisch{(-5.55,-4.75)}{fishblue}{0}{white}{fishyellow}{F4}{f}{0}
  \fisch{(-4.05,-4.75)}{fishblue}{1}{white}{fishred}{F8}{f}{0}
  \fisch{(-2.55,-4.75)}{fishblue}{1}{fishdark}{fishred}{F5}{h}{1}
  \fisch{(-1.05,-4.75)}{fishblue}{1}{fishdark}{fishyellow}{F6}{h}{1}
  \node[font=\footnotesize] at (-4.0,-6.13) {3 friedlich, 2 feindselig $\Rightarrow$ \emph{friedlich} $\Rightarrow$ \textbf{2 Fehler}};

  % Teilmenge Orange
  \draw[group,draw=fishorange,fill=white] (0.1,-3.45) rectangle (7.9,-6.5);
  \node[anchor=west,font=\small\bfseries,text=fishorange!80!black] at (0.3,-3.77) {Teilmenge Orange: 4 Fische};
  \fisch{(1.2,-4.75)}{fishorange}{0}{white}{fishyellow}{F1}{h}{0}
  \fisch{(3.0,-4.75)}{fishorange}{0}{white}{fishred}{F2}{h}{0}
  \fisch{(4.8,-4.75)}{fishorange}{1}{white}{fishyellow}{F7}{h}{0}
  \fisch{(6.6,-4.75)}{fishorange}{1}{fishdark}{fishred}{F9}{f}{1}
  \node[font=\footnotesize] at (4.0,-6.13) {1 friedlich, 3 feindselig $\Rightarrow$ \emph{feindselig} $\Rightarrow$ \textbf{1 Fehler}};

  % Ergebnis
  \node[draw=navy,fill=lightblue,rounded corners=1mm,inner sep=2mm,font=\small] at (0,-7.25)
    {Gesamtfehler nach dem Aufteilen: $2+1=3$ \qquad Informationsgewinn: $4-3=\mathbf{1}$};
\end{tikzpicture}

{\scriptsize\color{midgray}Körper = Schuppenfarbe, Punkte = Muster, unterer Streifen = Bauchfarbe, Flossen = Flossenfarbe.
Rot gestrichelt mit \glqq Fehler\grqq{}: Das Label des Fisches passt nicht zur Vorhersage seiner Gruppe.}
\end{center}

\smallskip
\begin{minipage}[t]{0.61\linewidth}
\vspace{0pt}
\renewcommand{\arraystretch}{1.1}\setlength{\tabcolsep}{4pt}
\begin{tabularx}{\linewidth}{@{}X>{\centering\arraybackslash}p{14mm}>{\centering\arraybackslash}p{16mm}>{\centering\arraybackslash}p{10mm}@{}}
\toprule
\textbf{Schuppenfarbe} & \textbf{friedlich} & \textbf{feindselig} & \textbf{Fehler}\\
\midrule
Fehler vor dem Aufteilen & & & 4\\
Blau & 3 & 2 & 2\\
Orange & 1 & 3 & 1\\
\midrule
Gesamtfehler nach dem Aufteilen & & & 3\\
\bottomrule
\end{tabularx}

\smallskip
Informationsgewinn: $4-3=1$
\end{minipage}\hfill
\begin{minipage}[t]{0.35\linewidth}
\vspace{0pt}
\renewcommand{\arraystretch}{1.1}\setlength{\tabcolsep}{4pt}
\begin{tabularx}{\linewidth}{@{}X>{\centering\arraybackslash}p{14mm}>{\centering\arraybackslash}p{6mm}@{}}
\toprule
\textbf{Attribut} & \textbf{Fehler}\newline\textbf{nachher} & \textbf{IG}\\
\midrule
\rowcolor{lightblue}\textbf{Schuppenfarbe} & 3 & \textbf{1}\\
Muster & 4 & 0\\
Bauchfarbe & 4 & 0\\
Flossenfarbe & 4 & 0\\
\bottomrule
\end{tabularx}

\smallskip
\emph{Schuppenfarbe} wird die Wurzel.
\end{minipage}
\end{beispiel}

\newpage
\sectionline{2. Wie entsteht ein Entscheidungsbaum?}

Für jede entstandene \fach{Teilmenge} wird erneut der Informationsgewinn aller Attribute bestimmt:
\begin{itemize}
  \item \textbf{Blau} (5 Fische, 2 Fehler): \emph{Muster} und \emph{Bauchfarbe} haben beide IG\,1. Bei diesem
        \fach{Gleichstand} sind beide gleich gut geeignet; im Beispiel wird \emph{Muster} gewählt.
  \item \textbf{Blau, Ohne} (2 Fische): beide friedlich, also ein \fach{Blatt}.
  \item \textbf{Blau, Punkte} (3 Fische, 1 Fehler) und \textbf{Orange} (4 Fische, 1 Fehler):
        \emph{Bauchfarbe} hat jeweils IG\,1 und trennt fehlerfrei.
\end{itemize}

\begin{center}
\begin{tikzpicture}[scale=0.82,transform shape,font=\small,every node/.style={align=center},
  decision/.style={draw=navy,fill=lightblue,rounded corners=1mm,minimum width=27mm,minimum height=9mm},
  leaf/.style={draw=midgray,fill=lightgray,rounded corners=1mm,minimum width=21mm,minimum height=9mm}]
  \node[decision] (root) at (0,0) {Schuppenfarbe?};
  \node[decision] (pattern) at (-4.2,-2.0) {Muster?};
  \node[decision] (bellyorange) at (4.6,-2.0) {Bauchfarbe?};
  \node[leaf] (none) at (-6.6,-4.0) {friedlich};
  \node[decision] (bellyblue) at (-1.8,-4.0) {Bauchfarbe?};
  \node[leaf] (bbBlack) at (-3.4,-6.0) {feindselig};
  \node[leaf] (bbWhite) at (-0.2,-6.0) {friedlich};
  \node[leaf] (obBlack) at (3.0,-4.0) {friedlich};
  \node[leaf] (obWhite) at (6.2,-4.0) {feindselig};
  \draw (root) -- node[above left,font=\footnotesize]{Blau} (pattern);
  \draw (root) -- node[above right,font=\footnotesize]{Orange} (bellyorange);
  \draw (pattern) -- node[above left,font=\footnotesize]{Ohne} (none);
  \draw (pattern) -- node[above right,font=\footnotesize]{Punkte} (bellyblue);
  \draw (bellyblue) -- node[above left,font=\footnotesize]{Schwarz} (bbBlack);
  \draw (bellyblue) -- node[above right,font=\footnotesize]{Weiß} (bbWhite);
  \draw (bellyorange) -- node[above left,font=\footnotesize]{Schwarz} (obBlack);
  \draw (bellyorange) -- node[above right,font=\footnotesize]{Weiß} (obWhite);
\end{tikzpicture}
\end{center}

\begin{beispiel}[title=Algorithmus zum Erstellen eines Entscheidungsbaums]
\begin{enumerate}
  \item Für die aktuelle Datenmenge den Informationsgewinn jedes Attributs bestimmen.
  \item Das Attribut mit dem größten Informationsgewinn wird zum \fach{Entscheidungsknoten}.
  \item Die Daten nach den Attributwerten in Teilmengen aufteilen.
  \item Haben alle Daten einer Teilmenge dasselbe Label oder ist keine sinnvolle Aufteilung mehr möglich,
        entsteht ein Blatt mit dem häufigeren Label. Sonst das Verfahren für diese Teilmenge wiederholen.
\end{enumerate}
\end{beispiel}

\begin{merke}
Ein Entscheidungsbaum entsteht schrittweise. In jedem Knoten wird ein möglichst geeignetes Attribut ausgewählt.
Anschließend wird das Verfahren für die entstandenen Teilmengen wiederholt.
\end{merke}

\end{document}
